\editionPage[\pdfbookmark[0]{Deux bilans annuels}{bm-bilans-guide-parent}]{bilans-guide}{Deux bilans annuels}{Sujets originaux · Réviser en reliant les notions}{Durée conseillée : 2 heures · Barème d’entraînement sur 20}
\begin{blocEdition}{bilans-guide}

Ces évaluations sont des créations du livre pour la Quatrième. Elles n’ont
pas de statut d’examen officiel. Le temps et le barème ci-dessous sont des
repères pour une séance autonome ; ils peuvent être adaptés en classe.
\editionEtape{Automatismes : environ 20 minutes}
Douze questions courtes, conseillées sans calculatrice : 6 points, soit
0,5 point par réponse. Écrire les étapes utiles, même lorsqu’elles sont brèves.
\editionEtape{Problèmes : environ 1 h 40}
Quatre problèmes indépendants, calculatrice autorisée : 14 points.
Chaque problème vaut 3,5 points : trois questions à 1 point et 0,5 point
pour la clarté de la rédaction, les unités et les contrôles.
\editionEtape{Travailler après l’évaluation}
Comparer les démarches, relever les erreurs par notion, puis refaire
les questions utiles sans correction. Un résultat faux peut accompagner
une démarche partiellement juste : identifier précisément l’étape à reprendre.
\editionEtape{Préparer une autre séance}
Revenir au cours et à deux exercices du chapitre concerné. Le second
bilan peut être fait plus tard pour vérifier que la méthode est devenue disponible.

\end{blocEdition}
\editionPage{bilan-A-auto}{Bilan A}{Calculer, lire et démontrer}{Partie 1 · 20 min · Sans calculatrice · 6 points}
\begin{blocEdition}{bilan-A-auto}
\begin{questions}
\item Calculer $(-8)\times(-3)$.
\item Calculer $7-2\times5$.
\item Calculer $\frac34+\frac18$.
\item Donner l’inverse de $-\frac25$.
\item Écrire $2^3\times2^4$ sous la forme d’une seule puissance.
\item Donner $\sqrt{121}$.
\item Développer $3(2x-5)$.
\item Résoudre $4x+3=19$.
\item Calculer 15\,\% de 80.
\item Par quel nombre multiplie-t-on un prix pour l’augmenter de 8\,\% ?
\item Donner la médiane de la série 2, 5, 8, 9, 12.
\item Une probabilité vaut 0,7. Donner celle de l’événement contraire.\end{questions}
\end{blocEdition}
\editionPage{bilan-A-ex-1}{Bilan A · Problème 1}{Deux ateliers}{Partie 2 · 1 h 40 pour les quatre problèmes · 3,5 points}
\begin{blocEdition}{bilan-A-ex-1}
Un atelier demande 12 euros d’inscription puis 4 euros par séance. Un autre atelier demande 7 euros par séance, sans inscription.\begin{questions}
\item Exprimer les deux dépenses $A(n)$ et $B(n)$ pour $n$ séances.
\item Pour combien de séances sont-elles égales ? Donner le prix commun.
\item Pour huit séances, on réduit de 10\,\% le prix total du premier atelier. Comparer les deux coûts.
\end{questions}\par\vspace{5mm}{\otLight\fontsize{8}{10}\selectfont\color{otMuted}Recherche et rédaction}\par\par\vspace{7mm}{\color{otRule!65}\hrule height .25pt}\par\vspace{7mm}{\color{otRule!65}\hrule height .25pt}\par\vspace{7mm}{\color{otRule!65}\hrule height .25pt}\par\vspace{7mm}{\color{otRule!65}\hrule height .25pt}\par\vspace{7mm}{\color{otRule!65}\hrule height .25pt}\par\vspace{7mm}{\color{otRule!65}\hrule height .25pt}
\end{blocEdition}
\editionPage{bilan-A-ex-2}{Bilan A · Problème 2}{Un terrain et un arroseur}{Partie 2 · 1 h 40 pour les quatre problèmes · 3,5 points}
\begin{blocEdition}{bilan-A-ex-2}
Le terrain ABC est un triangle rectangle en A, avec $AB=6$ m et $AC=8$ m.\begin{questions}
\item Calculer la longueur BC en justifiant.
\item On veut placer un arroseur à égale distance des trois coins A, B et C. Où le placer ? Quelle portée lui donner pour qu’il atteigne exactement les trois coins ?
\item Un câble suit $[BC]$ et nécessite 0,5 m de réserve à chaque extrémité. Il coûte 3,20 euros le mètre. Calculer le prix du câble.
\end{questions}\label{fig-bilan-A-ex-2}\figureGeo{bilan-a-geometrie}
\par\vspace{5mm}{\otLight\fontsize{8}{10}\selectfont\color{otMuted}Recherche et rédaction}\par\par\vspace{7mm}{\color{otRule!65}\hrule height .25pt}\par\vspace{7mm}{\color{otRule!65}\hrule height .25pt}\par\vspace{7mm}{\color{otRule!65}\hrule height .25pt}\par\vspace{7mm}{\color{otRule!65}\hrule height .25pt}\par\vspace{7mm}{\color{otRule!65}\hrule height .25pt}\par\vspace{7mm}{\color{otRule!65}\hrule height .25pt}
\end{blocEdition}
\editionPage{bilan-A-ex-3}{Bilan A · Problème 3}{Une petite enquête}{Partie 2 · 1 h 40 pour les quatre problèmes · 3,5 points}
\begin{blocEdition}{bilan-A-ex-3}
Douze élèves indiquent leur temps de lecture quotidien, en minutes : 5, 10, 10, 10, 15, 15, 20, 20, 20, 25, 25, 30.\begin{questions}
\item Calculer la moyenne au dixième, puis une médiane de la série.
\item Calculer l’étendue, puis la fréquence des durées d’au moins 20 minutes.
\item On choisit au hasard un de ces douze élèves. Calculer la probabilité que sa durée soit strictement inférieure à 20 minutes, puis vérifier avec l’événement contraire.
\end{questions}\par\vspace{5mm}{\otLight\fontsize{8}{10}\selectfont\color{otMuted}Recherche et rédaction}\par\par\vspace{7mm}{\color{otRule!65}\hrule height .25pt}\par\vspace{7mm}{\color{otRule!65}\hrule height .25pt}\par\vspace{7mm}{\color{otRule!65}\hrule height .25pt}\par\vspace{7mm}{\color{otRule!65}\hrule height .25pt}\par\vspace{7mm}{\color{otRule!65}\hrule height .25pt}\par\vspace{7mm}{\color{otRule!65}\hrule height .25pt}
\end{blocEdition}
\editionPage{bilan-A-ex-4}{Bilan A · Problème 4}{Un cône et un cylindre}{Partie 2 · 1 h 40 pour les quatre problèmes · 3,5 points}
\begin{blocEdition}{bilan-A-ex-4}
Un cône a pour rayon 3 cm et pour hauteur 12 cm.\begin{questions}
\item Calculer son volume exact, puis arrondi au dixième de cm$^3$.
\item Un cylindre a la même base et la même hauteur. Combien de cônes remplis faut-il verser pour remplir ce cylindre ?
\item Exprimer en mL la capacité du cône et dire si 110 mL de liquide peuvent y tenir.
\end{questions}\label{fig-bilan-A-ex-4}\figureGeo{bilan-a-cone}
\par\vspace{5mm}{\otLight\fontsize{8}{10}\selectfont\color{otMuted}Recherche et rédaction}\par\par\vspace{7mm}{\color{otRule!65}\hrule height .25pt}\par\vspace{7mm}{\color{otRule!65}\hrule height .25pt}\par\vspace{7mm}{\color{otRule!65}\hrule height .25pt}\par\vspace{7mm}{\color{otRule!65}\hrule height .25pt}\par\vspace{7mm}{\color{otRule!65}\hrule height .25pt}\par\vspace{7mm}{\color{otRule!65}\hrule height .25pt}
\end{blocEdition}
\editionPage{bilan-B-auto}{Bilan B}{Choisir une méthode et contrôler}{Partie 1 · 20 min · Sans calculatrice · 6 points}
\begin{blocEdition}{bilan-B-auto}
\begin{questions}
\item Calculer $(-45):9$.
\item Calculer $\frac25\times\frac{15}{4}$.
\item Simplifier $\frac{36}{48}$.
\item Écrire $5^3\times2^3$ sous la forme d’une puissance de 10.
\item Encadrer $\sqrt{90}$ entre deux entiers consécutifs.
\item Réduire $7x-4-2x+9$.
\item Factoriser $12x+18$.
\item Résoudre $5x=2x+12$.
\item Convertir 0,025 m$^3$ en litres.
\item Partager 35 dans le ratio $2:5$.
\item On lance deux pièces équilibrées. Quelle est la probabilité d’obtenir deux faces ?
\item Une translation déplace de 2 carreaux vers la gauche et de 3 vers le haut. Donner l’image du point $(4\,;1)$.\end{questions}
\end{blocEdition}
\editionPage{bilan-B-ex-1}{Bilan B · Problème 1}{Un bassin et deux unités}{Partie 2 · 1 h 40 pour les quatre problèmes · 3,5 points}
\begin{blocEdition}{bilan-B-ex-1}
Un bassin contient 720 L. Une pompe le vide à débit constant de $0{,}004$ m$^3$/s.\begin{questions}
\item Convertir le débit en L/s, puis en L/min.
\item Calculer la durée de vidange en minutes.
\item Après 1 min 30 s de fonctionnement, calculer le volume restant et la fraction de la capacité initiale qu’il représente.
\end{questions}\par\vspace{5mm}{\otLight\fontsize{8}{10}\selectfont\color{otMuted}Recherche et rédaction}\par\par\vspace{7mm}{\color{otRule!65}\hrule height .25pt}\par\vspace{7mm}{\color{otRule!65}\hrule height .25pt}\par\vspace{7mm}{\color{otRule!65}\hrule height .25pt}\par\vspace{7mm}{\color{otRule!65}\hrule height .25pt}\par\vspace{7mm}{\color{otRule!65}\hrule height .25pt}\par\vspace{7mm}{\color{otRule!65}\hrule height .25pt}
\end{blocEdition}
\editionPage{bilan-B-ex-2}{Bilan B · Problème 2}{Des milieux et un parallélogramme}{Partie 2 · 1 h 40 pour les quatre problèmes · 3,5 points}
\begin{blocEdition}{bilan-B-ex-2}
Le triangle ABC a pour côtés $AB=7$ cm, $AC=9$ cm et $BC=12$ cm. I, J et K sont les milieux de $[AB]$, $[AC]$ et $[BC]$.\begin{questions}
\item Démontrer que $(IJ)$ et $(BC)$ sont parallèles, puis calculer IJ.
\item Démontrer que le quadrilatère IBKJ est un parallélogramme.
\item Calculer le périmètre de IBKJ.
\end{questions}\label{fig-bilan-B-ex-2}\figureGeo{bilan-b-milieux}
\par\vspace{5mm}{\otLight\fontsize{8}{10}\selectfont\color{otMuted}Recherche et rédaction}\par\par\vspace{7mm}{\color{otRule!65}\hrule height .25pt}\par\vspace{7mm}{\color{otRule!65}\hrule height .25pt}\par\vspace{7mm}{\color{otRule!65}\hrule height .25pt}\par\vspace{7mm}{\color{otRule!65}\hrule height .25pt}\par\vspace{7mm}{\color{otRule!65}\hrule height .25pt}\par\vspace{7mm}{\color{otRule!65}\hrule height .25pt}
\end{blocEdition}
\editionPage{bilan-B-ex-3}{Bilan B · Problème 3}{Une série et un tirage}{Partie 2 · 1 h 40 pour les quatre problèmes · 3,5 points}
\begin{blocEdition}{bilan-B-ex-3}
Une boîte contient 20 cartes : six cartes portent le nombre 1, huit le nombre 2 et six le nombre 4. Les cartes sont indiscernables au toucher et l’on en tire une au hasard.\begin{questions}
\item Calculer la moyenne et la médiane des nombres inscrits.
\item Calculer la probabilité d’obtenir un nombre pair, puis celle de l’événement contraire.
\item On ajoute uniquement des cartes portant 1 pour que la probabilité d’un nombre pair devienne $\frac12$. Combien faut-il en ajouter ?
\end{questions}\par\vspace{5mm}{\otLight\fontsize{8}{10}\selectfont\color{otMuted}Recherche et rédaction}\par\par\vspace{7mm}{\color{otRule!65}\hrule height .25pt}\par\vspace{7mm}{\color{otRule!65}\hrule height .25pt}\par\vspace{7mm}{\color{otRule!65}\hrule height .25pt}\par\vspace{7mm}{\color{otRule!65}\hrule height .25pt}\par\vspace{7mm}{\color{otRule!65}\hrule height .25pt}\par\vspace{7mm}{\color{otRule!65}\hrule height .25pt}
\end{blocEdition}
\editionPage{bilan-B-ex-4}{Bilan B · Problème 4}{Tripler une variable}{Partie 2 · 1 h 40 pour les quatre problèmes · 3,5 points}
\begin{blocEdition}{bilan-B-ex-4}
Le script affiché utilise une seule variable x.\begin{questions}
\item Donner les valeurs de x après chacun des quatre tours.
\item Exprimer l’affichage avec une puissance, puis le calculer.
\item On remplace la valeur initiale 2 par une entrée $a$. Quelle entrée donne un affichage de 405 ?
\end{questions}\label{fig-bilan-B-ex-4}\figureGeo{bilan-b-script}
\par\vspace{5mm}{\otLight\fontsize{8}{10}\selectfont\color{otMuted}Recherche et rédaction}\par\par\vspace{7mm}{\color{otRule!65}\hrule height .25pt}\par\vspace{7mm}{\color{otRule!65}\hrule height .25pt}\par\vspace{7mm}{\color{otRule!65}\hrule height .25pt}\par\vspace{7mm}{\color{otRule!65}\hrule height .25pt}\par\vspace{7mm}{\color{otRule!65}\hrule height .25pt}\par\vspace{7mm}{\color{otRule!65}\hrule height .25pt}
\end{blocEdition}
